\uuid{6Mvb}
\exo7id{5103}
\titre{exo7 5103}
\auteur{rouget}
\organisation{exo7}
\datecreate{2010-06-30}
\isIndication{false}
\isCorrection{true}
\chapitre{Logique, ensemble, raisonnement}
\sousChapitre{Logique}
\cpge{mpsi-mp2i.raisonnement-vocabulaire-ensembliste.a, pcsi.raisonnement-vocabulaire-ensembliste.a, ptsi.raisonnement-vocabulaire-ensembliste.a, mpsi-mp2i.raisonnement-vocabulaire-ensembliste.c, pcsi.raisonnement-vocabulaire-ensembliste.d, ptsi.raisonnement-vocabulaire-ensembliste.d, bcpst1.logique-ensembles-raisonnement.a, bcpst1.vocabulaire-applications}
\module{Algèbre}
\niveau{L1}
\difficulte{}

\contenu{
\texte{
Exprimer à l'aide de quantificateurs les phrases suivantes puis donner leur négation.
}
\begin{enumerate}
    \item \question{($f$ étant une application du plan dans lui-même)
}
    \item \question{$f$ est l'identité du plan.}
\reponse{$(f=Id_{\mathcal{P}}\Leftrightarrow\forall M\in\mathcal{P},\;f(M)=M)$ et $(f\neq Id_{\mathcal{P}}\Leftrightarrow\exists
M\in\mathcal{P}/\;f(M)\neq M)$.}
    \item \question{$f$ a au moins un point invariant (on dit aussi point fixe).}
\reponse{$(f\;\mbox{a au moins un point fixe}\Leftrightarrow \exists M\in\mathcal{P}/\;f(M)=M)$ et $(f\;\mbox{n'a pas de point
fixe}\Leftrightarrow \forall M\in\mathcal{P},\;f(M)\neq M)$.}
\end{enumerate}
}
